\paragraph{A pole-free adjacent-order bound near a real zero.}
The spatial source coefficient appears singular if one
substitutes $-\partial_x^2\log|H_t(x)|$ directly at
a real zero. A simple algebraic cancellation removes
this artificial pole, and makes explicit the role
of \emph{other} real zeros.

\begin{proposition}[Nearest-root removal and a
pole-free Laguerre ratio]\label{prop:polefree}
Let $G$ be a nonconstant real entire function
in the Laguerre--P\'olya class of order at most
one, with infinitely many simple real zeros.
Fix a zero $a$. For $x$ in an open interval
containing $a$ but no other zeros, put
\[
S_a(x)=\sum_{\substack{\rho\in Z(G)\\\rho\ne a}}
\frac{1}{(x-\rho)^2}.
\]
Then $S_a$ is finite and continuous near $a$,
and for every integer $n\ge1$ and every such
real $x$,
\begin{equation}\label{eq:polefree}
\boxed{\quad
0\le
\frac{(n+1)(2n+1)L_{n+1}[G](x)}
{L_n[G](x)}
\le\frac{(n+1)(2n+1)}{n}S_a(x).
\quad}
\end{equation}
In particular $L_n[G](x)>0$ for all $n\ge1$,
and the bound holds \emph{at} $x=a$ without
introducing a singular factor $(x-a)^{-2}$.
\end{proposition}
\begin{proof}
For $x$ different from all zeros, the canonical
product gives
\[
\frac{G(x+iy)G(x-iy)}{G(x)^2}
=\prod_{\rho\in Z(G)}
\left(1+\frac{y^2}{(x-\rho)^2}\right).
\]
Write $b=(x-a)^{-2}$, and let $E_j$ be the
$j$th elementary symmetric function of
the positive numbers
$b_\rho=(x-\rho)^{-2}$ with $\rho\ne a$.
Then
\begin{align*}
L_n[G](x)&=G(x)^2(bE_{n-1}+E_n),\\
L_{n+1}[G](x)&=G(x)^2(bE_n+E_{n+1}).
\end{align*}
The finite elementary-symmetric inequality
$kE_k\le S_aE_{k-1}$ passes to the convergent
canonical-product coefficients by monotone
approximation. Thus
\[
nE_n\le S_aE_{n-1},\qquad
(n+1)E_{n+1}\le S_aE_n.
\]
For $b\ge0$,
\begin{align*}
n(bE_n+E_{n+1})
&\le S_a bE_{n-1}+S_a E_n\\
&=S_a(bE_{n-1}+E_n).
\end{align*}
Dividing by the positive denominator yields
\eqref{eq:polefree} away from $a$.
At $x=a$, factor $G(z)=(z-a)g(z)$.
The identity
\[
L_n[G](a)=G'(a)^2E_{n-1}(a)
\]
and continuity of all the other zero factors
give exactly the same inequality there.
The infinitude and simplicity of the real
zero set ensure the required coefficients
are positive at every finite order.
\end{proof}

The bound is conditional when applied to the
de Bruijn--Newman heat family: it is available
at those heat times where all the zeros are
known to be real and simple (notably $t=1/2$).
Extending it to the entire time slab
\emph{without assuming the unknown real-zero
property} is precisely the missing proof step.
There is no singularity from the selected
real zero; the genuine difficulty is controlling
the remaining inverse-square gap sum,
uniformly in heat time and Fourier height.

For example, a bound of the form
\[
S_a(t,x)\le C\bigl(1+\log^2(2+|x|)\bigr)
\]
for every heat time before a first collision
would imply a spatial coefficient of size at
most $C(n+1)(2n+1)n^{-1}
(1+\log^2(2+|x|))$ on the real-zero side.
Such a repulsion estimate is \emph{not} proved
for the genuine theta flow. It cannot be
inferred from average zero density alone,
and it must remain valid as the heat
parameter approaches any possible collision.

The elementary comparison is checked in
\path{XiKernel/PoleFreeLaguerre.lean} and by
exact rational examples in
\path{certificates/pole_free_laguerre_exact.py}.
The full infinite canonical-product identity
is an analytical input rather than an
end-to-end Lean theorem. The result is
a sharpened diagnostic for the remaining
heat-flow comparison, not a proof of RH.
